\documentclass[12pt, a4paper]{article}

% --- 常用宏包 ---
\usepackage[UTF8]{ctex}
\usepackage{amsmath, amsthm, amssymb, amsfonts}
\usepackage{geometry}
\usepackage{fancyhdr}
\usepackage{enumerate}
\usepackage{graphicx}
\usepackage{hyperref}
\usepackage{bm}
\usepackage{tcolorbox}

% --- 页面设置 ---
\geometry{left=2.5cm, right=2.5cm, top=2.5cm, bottom=2.5cm}
\pagestyle{fancy}
\fancyhf{}
\rhead{\today}
\lhead{近世代数作业}
\cfoot{\thepage}

% --- 环境定义 ---
\newtcolorbox{problem}[1]{
    colback=gray!5!white,
    colframe=gray!75!black,
    fonttitle=\bfseries,
    title=题目 #1
}

\newenvironment{solution}{
    \begin{proof}[\textbf{解}]
}{
    \end{proof}
}

% --- 个人信息 ---
\title{近世代数作业 05}
\author{李睿阳 (524120910059)}
\date{\today}

\begin{document}
\maketitle


\begin{problem}{1}
    设P是有限群G的Sylow p子群，证明：P是G的唯一Sylow p子群的充要条件是P为G的正规子群。
\end{problem}
\begin{solution}
    必要性：若 $P$ 是 $G$ 的唯一 Sylow $p$-子群，则对于 $\forall g \in G$，$gPg^{-1}$ 也是 $G$ 的 Sylow $p$-子群。由唯一性知 $gPg^{-1} = P$，故 $P \trianglelefteq G$。

    充分性：若 $P \trianglelefteq G$，则对于 $G$ 的任一 Sylow $p$-子群 $Q$，由 Sylow 第二定理知存在 $g \in G$ 使得 $Q = gPg^{-1}$。由于 $P$ 是正规子群，故 $Q = gPg^{-1} = P$，即 $P$ 是唯一的。
\end{solution}


\begin{problem}{2}
    设 $p,q$ 是不同素数，证明：任一 $pq$ 阶群都不是单群。
\end{problem}
\begin{solution}
    设 $|G| = pq$，不妨设 $p > q$。
    设 $n_p$ 为 $G$ 的 Sylow $p$-子群的个数。由 Sylow 第三定理：
    $n_p \equiv 1 \pmod p$ 且 $n_p \mid q$。
    由于 $q$ 是素数且 $q < p$，则 $n_p$ 只能为 $1$。
    由问题 1 知，该唯一的 Sylow $p$-子群 $P$ 是 $G$ 的正规子群。
    又 $|P| = p \neq 1, pq$，故 $P$ 是 $G$ 的非平凡正规子群，即 $G$ 不是单群。
\end{solution}


\begin{problem}{3}
    已知 $(a_1,\dots,a_n)$ 在 $G_1 \oplus \dots \oplus G_n$ 中的阶等于所有 $|a_i|$ 的最小公倍数。计算 $\mathbb{Z}_{45} \oplus \mathbb{Z}_{18}$ 中元素 $(40,12)$ 的阶。
\end{problem}
\begin{solution}
    $40$ 在 $\mathbb{Z}_{45}$ 中的阶为：$|40| = \frac{45}{\gcd(40,45)} = \frac{45}{5} = 9$。
    $12$ 在 $\mathbb{Z}_{18}$ 中的阶为：$|12| = \frac{18}{\gcd(12,18)} = \frac{18}{6} = 3$。
    故 $(40,12)$ 的阶为：$\operatorname{lcm}(9,3) = 9$。
\end{solution}
    

\begin{problem}{4}
    同构意义下列出所有 $1089$ 阶阿贝尔群。
\end{problem}
\begin{solution}
    $1089 = 3^2 \cdot 11^2$。
    由有限阿贝尔群基本定理
    \begin{enumerate}
        \item $\mathbb{Z}_{9} \oplus \mathbb{Z}_{121}$
        \item $\mathbb{Z}_{9} \oplus \mathbb{Z}_{11} \oplus \mathbb{Z}_{11}$
        \item $\mathbb{Z}_{3} \oplus \mathbb{Z}_{3} \oplus \mathbb{Z}_{121}$
        \item $\mathbb{Z}_{3} \oplus \mathbb{Z}_{3} \oplus \mathbb{Z}_{11} \oplus \mathbb{Z}_{11}$
    \end{enumerate}
\end{solution}
\end{document}
